% Focused spacing regression fixture for non-text elements.
% Compile from the repository root with:
%   lualatex -interaction=nonstopmode -halt-on-error -output-directory=tests tests/spacing-nontext-regression.tex
% Run at least three times after cleaning auxiliary files.
\DocumentMetadata{
  tagging=on,
  pdfversion=2.0,
  pdfstandard=ua-2,
  lang=en-US,
  testphase={phase-III,math}
}
\documentclass{ndsudisq}

\disquisitiontitle{Spacing Nontext Regression Fixture}
\disquisitionauthor{Template Test}
\documenttype{Dissertation}
\degree{Doctor of Philosophy}
\majordepartmentlabel{Major Department}
\program{Test Program}
\degreeoption{}
\finaldefense{May 2026}
\approvaldate{May 18, 2026}
\departmentchair{Department Chair}
\committeechair{Committee Chair}
\committeeone{Committee One}
\committeetwo{Committee Two}
\committeethree{Committee Three}
\subjectinfo{Fixture}
\keywordsinfo{NDSU, spacing, fixture}
\disquisitionlanguage{en-US}
\usecontinuouscaptions

\newcommand{\spacingfiller}{This ordinary body paragraph is intentionally long enough to wrap across multiple lines in the compiled PDF. It gives reviewers a stable double-spaced prose sample before or after a non-text element, making vertical spacing, paragraph indentation, and alignment changes easier to compare by eye.}

\begin{document}
\mainmattersetup

\chapter{Body Spacing and Nontext Regression}
\label{chap:spacing-body}

\spacingfiller

\section{Body Paragraphs Around Display Math}
\label{sec:body-display-math}
\spacingfiller

\mathalt{The energy E equals m c squared.}
\begin{equation}
E = mc^2
\label{eq:spacing-standard-display}
\end{equation}

\spacingfiller

A short lead-in line precedes this short display math sample:
\[
x+y=z
\]

\spacingfiller

\section{Body Lists Immediately After Headings}
\label{sec:body-lists}
\begin{itemize}
  \item This first body itemize item is deliberately long enough to wrap onto a second line so reviewers can compare line spacing within an item against the spacing between separate items in the same list.
  \item This second body itemize item provides a shorter comparison point after the multi-line item.
  \item This body item includes a nested list to exercise second-level list spacing.
    \begin{itemize}
      \item This nested item checks second-level itemize spacing after the parent item text.
      \item This second nested item checks spacing between adjacent second-level items.
    \end{itemize}
\end{itemize}

\spacingfiller

\subsection{Body Enumerate Immediately After a Heading}
\begin{enumerate}
  \item This first body enumerate item is deliberately long enough to wrap onto a second line so reviewers can compare wrapped item text, label alignment, and inter-item spacing.
  \item This second body enumerate item provides a shorter comparison point after the multi-line item.
\end{enumerate}

\spacingfiller

\subsection{Body Description Immediately After a Heading}
\begin{description}
  \item[Wrapped body term] This body description item is deliberately long enough to wrap onto a second line so reviewers can compare description-list looseness with the itemize and enumerate lists above.
  \item[Short body term] This second body description item provides a shorter comparison point after the multi-line description item.
\end{description}

\spacingfiller

\section{Floats, Captions, and Notes}
\label{sec:floats-captions-notes}
\spacingfiller

\begin{figure}[htbp]
\centering
\fbox{\rule{0pt}{1.5in}\rule{0.64\textwidth}{0pt}}
\caption{Short figure caption.}
\label{fig:spacing-short-caption}
\ndsufigurenote{Figure note text used to inspect note placement, note spacing, and note alignment relative to the short single-line caption above.}
\end{figure}

\spacingfiller

\begin{table}[htbp]
\centering
\caption{Longer wrapping table caption for visual comparison of caption alignment, indentation, and spacing when a title extends beyond a single line in the regression fixture.}
\label{tab:spacing-long-caption}
\begin{tabular}{lll}
\toprule
Column A & Column B & Column C \\
\midrule
Alpha & One & First value \\
Beta & Two & Second value \\
Gamma & Three & Third value \\
\bottomrule
\end{tabular}
\ndsutablenote{Table note text used to inspect note placement, note spacing, and note alignment relative to a longer wrapping caption and centered tabular material.}
\end{table}

\spacingfiller

\begin{schema}[htbp]
\centering
\fbox{\begin{minipage}{0.68\textwidth}
\centering
Schema placeholder: input $\rightarrow$ process $\rightarrow$ output.
\end{minipage}}
\caption{Schema caption with enough text to wrap and allow visual comparison with figure and table caption alignment in this fixture.}
\label{sch:spacing-schema}
\note{Schema note text. The class provides the generic \texttt{\textbackslash note} command and specialized figure/table note aliases; no schema-specific note alias is currently defined, so this fixture uses the recommended generic note command for schema notes.}
\end{schema}

\spacingfiller

\begin{figure}[p]
\centering
\fbox{\rule{0pt}{7.2in}\rule{0.78\textwidth}{0pt}}
\caption{Large float likely to occupy a float-only page and reveal whether float pages are top-aligned or vertically centered.}
\label{fig:spacing-float-page}
\ndsufigurenote{Large-float note for checking caption and note spacing on a likely float-only page.}
\end{figure}

\section{Code, Long Tables, and Boxed Placeholders}
\label{sec:other-nontext-spacing}
\spacingfiller

\begin{ndsucodelisting}[Template code listing spacing sample]
for row in rows:
    normalized = row.strip().lower()
    if normalized:
        print(normalized)
\end{ndsucodelisting}
\ndsucodenote{Code note used to inspect spacing after the template code-listing environment.}

\spacingfiller

\begin{center}
\fbox{\begin{minipage}{0.62\textwidth}
\centering
Boxed placeholder outside a float.\\
This simulates figure-like content used by lightweight tests.
\end{minipage}}
\end{center}

\spacingfiller

\begin{longtable}{lll}
\caption{Longtable caption included to check longtable caption spacing and alignment.}\label{tab:spacing-longtable}\\
\toprule
Entry & Category & Note \\
\midrule
\endfirsthead
\caption[]{Longtable caption included to check longtable caption spacing and alignment, continued.}\\
\toprule
Entry & Category & Note \\
\midrule
\endhead
\midrule
\multicolumn{3}{r}{Continued on next page}\\
\endfoot
\bottomrule
\endlastfoot
One & A & This row begins the longtable spacing sample. \\
Two & A & This row provides additional table body rhythm. \\
Three & B & This row provides additional table body rhythm. \\
Four & B & This row provides additional table body rhythm. \\
Five & C & This row provides additional table body rhythm. \\
Six & C & This row provides additional table body rhythm. \\
Seven & D & This row provides additional table body rhythm. \\
Eight & D & This row provides additional table body rhythm. \\
Nine & E & This row provides additional table body rhythm. \\
Ten & E & This row ends the compact longtable sample. \\
\end{longtable}

\spacingfiller

\clearpage
\chapter{Heading After a Forced Page Break}
\label{chap:forced-break-heading}
\spacingfiller

\section{Paragraphs After Forced-Break Chapter Heading}
\spacingfiller

\appendix
\chapter{Appendix List Spacing Regression}
\label{app:list-spacing}

\section{Appendix Itemize Immediately After a Heading}
\begin{itemize}
  \item This first appendix itemize item is deliberately long enough to wrap onto a second line so reviewers can compare appendix itemize spacing against body itemize spacing.
  \item This second appendix itemize item provides a shorter comparison point after the multi-line item.
\end{itemize}

\spacingfiller

\section{Appendix Enumerate Immediately After a Heading}
\begin{enumerate}
  \item This first appendix enumerate item is deliberately long enough to wrap onto a second line so reviewers can compare appendix enumerate spacing against body enumerate spacing.
  \item This second appendix enumerate item provides a shorter comparison point after the multi-line item.
\end{enumerate}

\spacingfiller

\section{Appendix Description Immediately After a Heading}
\begin{description}
  \item[Wrapped appendix term] This appendix description item is deliberately long enough to wrap onto a second line so reviewers can compare description-list looseness against appendix itemize and enumerate lists.
  \item[Short appendix term] This second appendix description item provides a shorter comparison point after the multi-line appendix description item.
\end{description}

\spacingfiller

\end{document}
